% Fragmento integrable. Preambulo anfitrion: amsmath, amssymb, longtable, hyperref.
\section{Ordered transport and Jacobi representation}
\label{hmtfg:fragmento:transporte-ordenado-y-representacion-de-jacobi}

\subsection{Complete transport from the continuum and spectral realization of the profile}
\label{hmtfg:escalado:14}

The ordered transport of the continuum and the spectral representation of the reader are composed before parameter selection.

\subsubsection{Coverage of all roots and transport of their minimum}
\label{hmtfg:escalado:14.1}

The ordered construction of the continuum developed in the preceding works produces the surjective Archimedean readout from compatible cylinders and their ordered quotient. Within its declared universe \(\mathfrak G\), it proves the isomorphism of complete ordered fields
\begin{equation}
\iota:\mathbb R_{\rm HMT}^{\mathfrak G}
   \xrightarrow{\cong}\mathbb R^{\mathfrak G}.
\label{hmtfg:escalado:eq:37}
\end{equation}
The name \(\eta\) used in the source work is changed here to \(\iota\), to distinguish it from the weighting parameter, fixed at zero. The second projection preserves incidence, boundary and memory over the same antecedent; the evaluation quotient and the complete history retain their respective domains.

Let us construct in the HMT field, before selecting a root,
\[
\cos_{\rm H}(t)=\sum_{k=0}^{\infty}\frac{(-1)^kt^{2k}}{(2k)!},
\qquad
u_{\rm H}(x)=\frac1{12}\sum_{m=1}^{12}
\left[1-\cos_{\rm H}\!\left(\frac{2(3m-1)x}{\sqrt[\rm H]{899}}\right)\right],
\]
\begin{equation}
F_a(x)=1-a\,u_{\rm H}(x),\qquad S_n^{\rm H}(a)=F_a^{\,2^n}(0).
\label{hmtfg:escalado:eq:38}
\end{equation}
The integers and the second moment are those already derived from the dodecaphase reader. The square root is the positive one in the ordered field. The convergent series analytically realizes that reader; it preserves the profile and the normalizing coefficient of section~\ref{hmtfg:escalado:1}.

\textbf{Return transport theorem.} For every parameter in the internal domain of \eqref{hmtfg:escalado:eq:38},
\begin{equation}
\iota(F_a(x))=f_{\iota(a)}(\iota(x)),\qquad
\iota(S_n^{\rm H}(a))=S_n(\iota(a)).
\label{hmtfg:escalado:eq:39}
\end{equation}
\textbf{Proof.} The ordered isomorphism fixes the rationals and preserves the positive square root. It preserves convergent limits: every neighborhood of positive rational radius is transported to another of the same radius. Applied to the partial sums of \eqref{hmtfg:escalado:eq:38}, it transports the cosine and \(u_{\rm H}\). The first equality follows from preservation of addition and multiplication; the second follows by induction on the number of iterations.

With a previous center \(a_{n-1}\) fixed, define \(Z_n^{\rm H}\) by zero return, exact critical period \(2^n\) and \(a>a_{n-1}\), within the declared parameter domain. Let \(Z_n\) be the set defined by the same conditions for \(f_\lambda\), to the right of \(\iota(a_{n-1})\). Then
\begin{equation}
\boxed{\iota[Z_n^{\rm H}]=Z_n,\qquad
\iota(\min Z_n^{\rm H})=\min Z_n.}
\label{hmtfg:escalado:eq:40}
\end{equation}
The second equality is asserted when the minimum exists; the first also preserves the nonexistence of candidates.

\textbf{Proof.} \eqref{hmtfg:escalado:eq:39} preserves the zero return and the earlier returns determining the exact period. Order preserves position relative to the previous center. Surjectivity of \eqref{hmtfg:escalado:eq:37}, applied to any root parameter, proves reverse coverage. If \(a_*\) is the minimum, every \(z\in Z_n\) has the form \(\iota(a)\), with \(a\in Z_n^{\rm H}\), whence \(\iota(a_*)\le z\). The converse assertion uses \(\iota^{-1}\).

Surjectivity covers every parameter in the chart, including those whose roots have not been isolated by a finite computation. The isomorphism preserves the declared universe of interpretation and the selection order.

\subsubsection{Jacobi representation of the complete dodecaphase profile}
\label{hmtfg:escalado:14.2}

Define the atomic measure of the reader,
\begin{equation}
s_m=\frac{4(3m-1)^2}{899},\qquad
\nu_0=\frac1{12}\sum_{m=1}^{12}\delta_{s_m},\qquad
M_r=\int s^r\,d\nu_0(s).
\label{hmtfg:escalado:eq:41}
\end{equation}
The twelve nodes are distinct and positive. For a polynomial \(p\ne0\) of degree less than \(r\le12\),
\begin{equation}
\sum_{i,j=0}^{r-1}p_iM_{i+j}p_j
=\frac1{12}\sum_{m=1}^{12}p(s_m)^2>0.
\label{hmtfg:escalado:eq:42}
\end{equation}
A polynomial of degree less than twelve does not vanish at all twelve nodes. At degree twelve, \(\prod_m(s-s_m)\) appears, with zero norm.

Orthogonalization with respect to this measure produces orthonormal polynomials \(\psi_0,\ldots,\psi_{11}\), with \(\psi_0=1\). Let
\[
Q_{mj}=\psi_j(s_m)/\sqrt{12},\quad
D=\operatorname{diag}(s_m),\quad J_{12}=Q^{\mathsf T}DQ.
\]
The following hold
\begin{equation}
Q^{\mathsf T}Q=I,\quad DQ=QJ_{12},\qquad
\boxed{u_0(x)=e_0^{\mathsf T}
\bigl[I-\cos(x\sqrt{J_{12}})\bigr]e_0.}
\label{hmtfg:escalado:eq:43}
\end{equation}
\textbf{Proof.} Orthogonality follows from \eqref{hmtfg:escalado:eq:42}. Multiplication by \(s\) produces a three-term recurrence: entries separated by more than one position vanish by orthogonality and degree. Choosing positive leading coefficients makes the subdiagonals positive. By functional calculus,
\(\cos(x\sqrt{J_{12}})=Q^{\mathsf T}\cos(x\sqrt D)Q\).
The vector \(Qe_0\) is constant, with value \(1/\sqrt{12}\); substituting it recovers the sum in section~\ref{hmtfg:escalado:1}.

The article on Grassmannians develops this mechanism for marked measures and refinements. Here it is applied to \eqref{hmtfg:escalado:eq:41}; rank twelve belongs to this reader. The genealogical marks are retained alongside the reader, in accordance with the marked faithfulness of the source work. \(J_{12}\) represents the analytic profile; the complete state additionally preserves routes, orientation and memory.

The verifier for this extension checks, with exact fractions, the twelve degrees of positive norm, vanishing at degree twelve, the intertwiner \(VT=DV\) in the monic basis and weighted self-adjointness \(HT=T^{\mathsf T}H\). It checks moments of degree zero through thirty; the identity for every degree follows from the intertwiner. Its first coefficients are
\begin{equation}
a_0=2,\qquad \beta_1=\frac{2493348}{808201}.
\label{hmtfg:escalado:eq:44}
\end{equation}

\subsubsection{Oriented sensitivity of the return}
\label{hmtfg:escalado:14.3}

For an exact return of period \(p=2^n\), let
\[
c_j=f_\lambda^j(0),\quad q_j=\partial_\lambda c_j,\quad
D_j=\prod_{k=1}^j f_\lambda'(c_k),\quad D_0=1.
\]
Before the critical return, the factors of \(D_{p-1}\) are nonzero. Differentiation and telescoping give
\begin{equation}
q_{j+1}=-u_0(c_j)+f_\lambda'(c_j)q_j,\qquad
\boxed{\frac{q_p}{D_{p-1}}
=-\sum_{j=1}^{p-1}\frac{u_0(c_j)}{D_j}.}
\label{hmtfg:escalado:eq:45}
\end{equation}
In the canonical word, \(\operatorname{sign}c_j=(-1)^{v_2(j)}\). The sign of \(f_\lambda'(c_j)\) is opposite; therefore,
\begin{equation}
\operatorname{sign}D_j=(-1)^{j+\sum_{k=1}^jv_2(k)}
=(-1)^{s_2(j)},
\label{hmtfg:escalado:eq:46}
\end{equation}
using \(\sum_{k=1}^jv_2(k)=j-s_2(j)\), where \(s_2(j)\) counts the binary ones. The normalized sensitivity is
\begin{equation}
\mathcal T_n=-\sum_{j=1}^{2^n-1}(-1)^{s_2(j)}t_j,\qquad
t_j=\frac{u_0(c_j)}{|D_j|}>0.
\label{hmtfg:escalado:eq:47}
\end{equation}
Nodal positivity of \eqref{hmtfg:escalado:eq:42} determines \(u_0\); \eqref{hmtfg:escalado:eq:47} additionally retains the orbital products and their orientations.

The additional moment representation would require a positive measure \(\rho\) on \([0,1]\) with \(t_j=\int z^{j-1}\,d\rho\). Under that condition,
\begin{equation}
\mathcal T_n=\int_0^1
\frac{1-\prod_{r=0}^{n-1}(1-z^{2^r})}{z}\,d\rho(z)>0
\label{hmtfg:escalado:eq:48}
\end{equation}
for \(\rho\ne0\), extending the integrand by its value \(1\) at zero. The identity is obtained by expanding the binary product.

The rational check at the previously certified period-four root obtains
\[
t_1\approx0.40425224267600139730,\quad
t_2\approx0.04925249167432646403,\quad
t_3\approx0.15740870098294833737
\]
and the rigorous outer interval
\begin{equation}
0.296096033367379523967764444238832345360107
<\mathcal T_2<
0.296096033367379523967764444238832345360112.
\label{hmtfg:escalado:eq:49}
\end{equation}
It also certifies \(t_3>t_2\). A positive measure on \([0,1]\) would impose \(t_3\le t_2\); that specific lifting of orbital weights to positive moments is ruled out. The positive sign of \eqref{hmtfg:escalado:eq:49} remains certified. This check distinguishes the failure of a proposed proof from the failure of the property it sought to prove.
